\section{Major theorems}
\label{sec:theorems}

The library proves the central existence and structural theorems
of finite non-cooperative game theory.  The proofs are stated on
the kernel-game/game-form layer (or on layered information-model
abstractions; see \S\ref{sec:theorems:kuhn}), and inherited by
every language via $\compileTo$.  This section sketches their
content; full proofs are in the cited Lean files.

\subsection{Existence of mixed Nash equilibria}
\label{sec:theorems:nash-mixed}

\begin{theorem}[Mixed Nash existence; \TT{KernelGame.mixed\_nash\_exists}]
\label{thm:nash-mixed}
Let $\KG$ be a kernel game with $\Players$ finite, every $\Strat_i$
finite and non-empty, and $\Outcome$ finite.  Then the mixed
extension $\KG^{\mathrm{mix}}$ admits a Nash equilibrium:\ a
profile $\Profile^* \in \prod_i \PMF{\Strat_i}$ with
\[
  \forall i,\,\forall \tau \in \PMF{\Strat_i}.\quad
    \EU(\KG^{\mathrm{mix}}, \Profile^*, i)
      \;\ge\;
    \EU(\KG^{\mathrm{mix}}, \Profile^*[i \mapsto \tau], i).
\]
\end{theorem}

\noindent\emph{Proof structure.}\ Following Nash's
original~\citep{Nash1950, Nash1951}, we define the gain map
\[
  \Phi_i(\Profile)(s)
    \;=\; \frac{\Profile_i(s) + (\mathrm{gain}_i(\Profile, s))^+}
              {1 + \sum_{s'} (\mathrm{gain}_i(\Profile, s'))^+},
  \qquad
  \mathrm{gain}_i(\Profile, s)
    \;=\;
    \EU(\KG^{\mathrm{mix}}, \Profile[i \mapsto \Dirac_s], i)
      - \EU(\KG^{\mathrm{mix}}, \Profile, i),
\]
which sends the product simplex $\prod_i \Delta(\Strat_i)$
continuously into itself.  The Brouwer fixed-point theorem from
the \TT{fixed-point-theorems-lean4}
package~\citep{harfeFixedPointTheorems} is lifted to the product
of standard simplices, and a fixed point of $\Phi$ is shown to be
a Nash equilibrium by an algebraic argument bracketing the positive
parts.

\paragraph{Why mathlib's Brouwer.}  The continuity content of the
proof — that the Nash map is continuous on a compact convex
polytope — is genuinely topological, and mechanizing the topology
of the simplex from scratch would dwarf the rest of the proof.
mathlib already provides a Brouwer-on-the-disk theorem and the
infrastructure to lift it across homeomorphisms; we use the
standard simplex's homeomorphism to a closed ball.  The product
of simplices step is handled by an explicit product-fixed-point
lemma rather than by appealing to the product theorem in topology.

\tradeoff{Brouwer is non-constructive, so mixed Nash existence is
\TT{noncomputable}.  A constructive existence proof would require
either an explicit algorithmic construction
(Lemke–Howson~\citep{LemkeHowson1964} for two-player games, or
the support-enumeration class for finite games more generally) or
a constructive fixed-point theorem on a finite simplicial
structure (Sperner's lemma~\citep{Sperner1928} suffices for
two-player games but the multi-player extension is not as clean).
We took the topological route because an existing Lean~4
mechanization of the Brouwer fixed-point theorem was available
as the standalone \TT{fixed-point-theorems-lean4}
package~\citep{harfeFixedPointTheorems}, which proves Brouwer for
nonempty compact convex subsets of finite-dimensional normed
spaces via a cubical Sperner's-lemma route (Kuhn 1960); this
gave us a ready dependency in the right shape for arbitrary $n$
players and let us focus on the strategic-form content of the
proof.  Algorithm-level Nash existence is independent enough to
live in a downstream module.}

\leancorr{Statement \TT{KernelGame.mixed\_nash\_exists} in
\TT{Theorems/NashExistenceMixed.lean}; the product-simplex
Brouwer wrapper in \TT{Concepts/ProductSimplexBrouwer.lean}.
Supplement~§6.1.}

\subsection{Existence of correlated and coarse correlated equilibria}
\label{sec:theorems:correlated-existence}

\begin{theorem}[CE / CCE existence; \TT{KernelGame.correlatedEq\_exists}, \TT{coarseCorrelatedEq\_exists}]
\label{thm:ce-exists}
Under the same finiteness hypotheses as
\Cref{thm:nash-mixed}, $\KG$ admits a correlated equilibrium and a
coarse correlated equilibrium.
\end{theorem}

\noindent\emph{Proof.}\ The independent product distribution
$\bigotimes_i \Profile^*_i$ over the mixed Nash profile
$\Profile^*$ of \Cref{thm:nash-mixed} is a correlated equilibrium
of the original $\KG$ (every recommendation-dependent deviation
factors through a constant-strategy unilateral deviation in
$\KG^{\mathrm{mix}}$).  CCE existence is the corollary
$\Pref\text{-CE}\subseteq\Pref\text{-CCE}$
(\S\ref{sec:solution-concepts:correlated}).

\tradeoff{This route gives CE existence as a corollary of
mixed Nash existence and is the cleanest path through the
finite-game machinery already in the library.  Hart and
Mas-Colell's regret-matching no-regret learning
construction~\citep{HartMasColell2000} produces CCE existence
constructively without going through Brouwer.  We have not
formalized regret-matching itself, but the existence theorem is
still available as a corollary of Nash existence; if a future
constructive Nash-existence proof becomes available the dependency
inverts trivially.}

\leancorr{\TT{KernelGame.correlatedEq\_exists} and
\TT{coarseCorrelatedEq\_exists} in
\TT{Theorems/CorrelatedEqExistence.lean}.  Supplement~§6.2.}

\subsection{The minimax theorem}
\label{sec:theorems:minimax}

\begin{theorem}[Von Neumann's minimax; \TT{KernelGame.von\_neumann\_minimax}]
\label{thm:minimax}
Let $\KG$ be a finite two-player zero-sum kernel game.  There exist
a real value $v$ and a mixed Nash profile $\sigma$ of
$\KG^{\mathrm{mix}}$ such that player $1$ cannot push player $0$'s
payoff below $v$ by deviating from $\sigma_1$, and player $0$ cannot
push their payoff above $v$ by deviating from $\sigma_0$:
\[
  \EU(\KG^{\mathrm{mix}}, \sigma, 0) = v,\qquad
  \forall s_1,\ \EU(\KG^{\mathrm{mix}}, \sigma[1 \mapsto s_1], 0) \ge v,
\]
\[
  \forall s_0,\ \EU(\KG^{\mathrm{mix}}, \sigma[0 \mapsto s_0], 0) \le v.
\]
Moreover, all mixed Nash equilibria of a two-player zero-sum game
give player $0$ the same expected utility.
\end{theorem}

\noindent\emph{Proof.}\ \Cref{thm:nash-mixed} gives existence of a
mixed Nash equilibrium.  In a two-player zero-sum game, mixed Nash
profiles give saddle-style inequalities for $\EU(\cdot, 0)$,
which combine with the zero-sum-Nash equality lemma and mixed Nash
existence to deliver the displayed guarantees.  The classical
max-min/min-max equality is the standard mathematical reading of
the finite saddle inequalities; the Lean theorem is stated in the
guarantee form above.

\leancorr{Statement \TT{KernelGame.von\_neumann\_minimax} in
\TT{Theorems/Minimax.lean}; supporting zero-sum-Nash equality in
\TT{Concepts/ZeroSumNash.lean} (\TT{IsZeroSum.nash\_eu\_eq}).
Supplement~§6.3.}

\subsection{Backward induction (Zermelo) and the one-shot deviation principle}
\label{sec:theorems:zermelo-odp}

A \emph{subgame-perfect equilibrium}~\citep{Selten1965} (SPE) of
an extensive-form game is a strategy profile that is a Nash
equilibrium not just of the whole game but of every subgame:\ the
restriction of $\sigma$ to any decision node and its descendants
is itself Nash on the subtree rooted there.  This rules out
``empty threats''  — Nash profiles that prescribe a player to do
something off-path that they would not actually want to do if the
off-path subgame were reached.  In a finite perfect-information
tree, SPE is the standard refinement of Nash equilibrium, and is
the equilibrium concept the theorems below characterize.

\begin{theorem}[One-shot deviation principle]
\label{thm:odp}
Let $G$ be a finite perfect-information extensive-form game.  A pure
profile $\sigma$ is a subgame-perfect equilibrium of $G$ iff it has
no profitable one-shot deviation:\ at every reachable decision node,
the action prescribed by $\sigma$ yields at least as much expected
utility (holding the rest of $\sigma$ fixed) as any alternative.
\end{theorem}

\begin{theorem}[Zermelo's backward induction]
\label{thm:zermelo}
Every finite perfect-information extensive-form game has a pure
subgame-perfect equilibrium.
\end{theorem}

\noindent The backward induction itself constructs, by recursion
on the tree, a profile with no profitable one-shot deviation; the
ODP closes the loop to subgame-perfection.  Two non-trivial
ingredients are worth flagging:
\begin{itemize}\itemsep2pt
\item \emph{Disjointness of subtrees in perfect-info trees}:\ in
a perfect-info tree, sibling subtrees have disjoint information-set
supports.  This is what allows backward induction to fix actions
at a decision node without breaking optimality elsewhere.

\item \emph{Reachable-decision quantification}:\ the ODP's
hypothesis quantifies over reachable decision nodes only.  An
explicit reach-by relation makes this precise and lets the proof
ignore unreachable nodes whose action choice cannot affect EU.
\end{itemize}

\tradeoff{Backward induction is constructively well-defined on a
finite tree, but the library's Zermelo is \TT{noncomputable}
because intermediate steps invoke classical reasoning to compare
real-valued utilities.  A computable variant restricting to
rational utilities would compile but is not packaged.}

\leancorr{Statements \TT{EFG.oneShotDeviation\_iff\_spe} and
\TT{EFG.zermelo} in \TT{Theorems/OneShotDeviation.lean} and
\TT{Theorems/Zermelo.lean}; the subtree disjointness lemma is
\TT{perfectInfo\_disjoint\_subtrees}; the reachability relation
is \TT{EFG.ReachBy}.  Supplement~§6.4.}

\subsection{Kuhn's theorem: \texorpdfstring{behavioral $\leftrightarrow$ mixed strategies}{behavioral vs. mixed strategies}}
\label{sec:theorems:kuhn}

Kuhn's theorem~\citep{Kuhn1953} relates two carriers of randomness
in extensive-form games:\ a \emph{mixed strategy} is a
distribution over pure strategies (a single global lottery),
while a \emph{behavioral strategy} chooses an independent local
lottery at each information set.  The two are payoff-equivalent
under recall conditions; without recall, the asymmetry is
genuine.

The library proves Kuhn's theorem on a representation-neutral
substrate, the \emph{observation model}:\ a tuple
$(\sigma, \mathrm{Obs}, \Act)$ recording the world state, each
player's observation type, and the per-observation action arity.
A behavioral profile is a per-player function
$\InfoState_i \to \PMF{\Act_i}$; a mixed profile is a per-player
distribution over the deterministic functions
$\InfoState_i \to \Act_i$.\footnote{Notational reconciliation:\
the Lean code uses \TT{runDist}/\TT{runDistPure} for the
profile-evaluation function on the observation-model carrier; we
write $\evalDist$ throughout the chapter for uniformity with the
EFG-language naming.  The two refer to the same operation
specialized to the relevant carrier.}

\begin{theorem}[Kuhn behavioral $\to$ mixed]
\label{thm:kuhn-bm}
Under the bounded-horizon hypothesis (no infinite traces) and the
\emph{no-non-trivial info-state repeat} hypothesis, every
behavioral profile $\beta$ admits an equivalent mixed profile
$\mu(\beta)$ in the sense
\[
  \evalDist(\beta) \;=\; \mu(\beta) \mathbin{\bind} \evalPure.
\]
\end{theorem}

\begin{theorem}[Kuhn mixed $\to$ behavioral]
\label{thm:kuhn-mb}
Under per-step semantic locality assumptions (sufficient for the
\emph{ActionPosteriorLocal} property), every mixed profile $\mu$
admits an equivalent behavioral profile $\beta(\mu)$ in the sense
\[
  \evalDist(\beta(\mu)) \;=\; \mu \mathbin{\bind} \evalPure.
\]
\end{theorem}

\noindent The two directions are not symmetric.  $B \to M$ does
not need recall — the mixed lottery can be constructed
ex-ante by independent product across information states — but
does need horizon-separation (which holds on snapshot-refined
$\InfoState$'s).  $M \to B$ requires the per-step recall property
that, given a player's information state, the conditional
distribution of their action under $\mu$ does not depend on
information \emph{outside} the player's recall.  This is
strictly weaker than perfect recall in the textbook sense:\ it
does not require remembering one's own past observations
literally, only that the action conditional has the right
locality.

The snapshot-refined information-state structure is packaged as a
derived layer on top of an arbitrary observation model.  A second
proof of the $M \to B$ direction is given via the mixed-Nash-induced
correlated distribution.  Language-specific Kuhn theorems for EFG,
MAID, MR, and FOSG follow by composing the language's compilation
to the observation-model layer with the generic theorem.

\tradeoff{Stating Kuhn on the observation-model layer rather than
on a specific syntactic class of EFGs is the single largest piece
of abstraction work in the library.  The benefit is that one proof
serves every language with a generic compilation through that
layer; the cost is a stack of technical infrastructure
(information-state snapshots, the locality predicates) that is not
in the textbook proof.  We judge the tradeoff to be sharply
positive:\ language-specific Kuhn theorems become single-line
corollaries of a single core result.}

\leancorr{\TT{Theorems/Kuhn/} contains \TT{ObsModel} (the
snapshot-refined wrapper), \TT{ObsModelCore}, \TT{KuhnModel}, the
two proof files \TT{BehavioralToMixedCore} and
\TT{MixedToBehavioralCore}, and \TT{CorrelatedRealization} (the
alternative $M \to B$ route).  The hypothesis predicates
\TT{NoNontrivialInfoStateRepeat}, \TT{HorizonSeparation},
\TT{ActionPosteriorLocal}, \TT{PerStepPlayerRecall} live on
\TT{ObsModelCore}.  Supplement~§6.5.}

\subsection{Welfare and welfare theorems}
\label{sec:theorems:welfare}

The library packages the basic welfare analysis:\ social welfare
is well-defined under finiteness, every
team-game welfare maximizer is Nash, and the price of stability
in a team game is $1$.  We do not formalize the deeper welfare
theorems of mechanism design here — they live in
\S\ref{sec:mechanism-auctions}.

\leancorr{\TT{Concepts/WelfareTheorems.lean}.}
